%% ============================================================
%%  SECCIÓN 6 — Banco de Problemas Resueltos y Ejercicios Propuestos
%% ============================================================

\section{Banco de Problemas Resueltos y Ejercicios Propuestos}
\label{sec:ejercicios}

\subsection{Problemas Resueltos Paso a Paso}

\subsubsection[Problema 1: Rotaciones finitas vía Hadamard]{Problema Resuelto 1: Rotaciones Espaciales Cuánticas Finitas vía Hadamard}
\begin{ejemplobox}[Problema Resuelto 1: Rotaciones Espaciales Cuánticas Finitas vía Hadamard]
    \textbf{Enunciado:} El operador de rotación finita por un ángulo $\theta \in \mathbb{R}$ alrededor del eje polar $z$ está dado por:
    \begin{equation}
        \hat{R}_z(\theta) \equiv \exp\left(-\frac{\iu}{\hbar}\theta\hat{L}_z\right).
    \end{equation}
    \begin{enumerate}
        \item Calcule los conmutadores fundamentales $[\hat{L}_z, \hat{x}]$ y $[\hat{L}_z, \hat{y}]$.
        \item Utilice el Lema de Hadamard para deducir la forma del operador transformado $\hat{x}' \equiv \hat{R}_z\adj(\theta)\hat{x}\hat{R}_z(\theta)$ y verifique que reproduce la matriz de rotación geométrica estándar en el plano $xy$:
        \begin{equation}
            \hat{x}' = \hat{x}\cos\theta - \hat{y}\sin\theta.
        \end{equation}
    \end{enumerate}

    \textbf{Solución Paso a Paso:}

    \textit{Paso 1: Cálculo de Conmutadores de Posición con $\hat{L}_z$:}
    Escribimos $\hat{L}_z = \hat{x}\hat{p}_y - \hat{y}\hat{p}_x$. Evaluamos el conmutador con la coordenada $\hat{x}$:
    \begin{align}
        [\hat{L}_z, \hat{x}] &= [\hat{x}\hat{p}_y - \hat{y}\hat{p}_x, \, \hat{x}] = [\hat{x}\hat{p}_y, \hat{x}] - [\hat{y}\hat{p}_x, \hat{x}].
    \end{align}
    Como $\hat{x}$ conmuta con $\hat{p}_y$ y consigo mismo, $[\hat{x}\hat{p}_y, \hat{x}] = 0$. Por ende:
    \begin{equation}
        [\hat{L}_z, \hat{x}] = -[\hat{y}\hat{p}_x, \hat{x}] = -\hat{y}[\hat{p}_x, \hat{x}] - [\hat{y}, \hat{x}]\hat{p}_x.
    \end{equation}
    Sabiendo que $[\hat{p}_x, \hat{x}] = -\iu\hbar\hat{I}$ y $[\hat{y}, \hat{x}] = 0$:
    \begin{equation}
        [\hat{L}_z, \hat{x}] = -\hat{y}(-\iu\hbar\hat{I}) = \boxed{\iu\hbar\hat{y}}.
        \label{eq:lz_x_comm}
    \end{equation}
    Análogamente, para la coordenada $\hat{y}$:
    \begin{equation}
        [\hat{L}_z, \hat{y}] = [\hat{x}\hat{p}_y - \hat{y}\hat{p}_x, \, \hat{y}] = [\hat{x}\hat{p}_y, \hat{y}] = \hat{x}[\hat{p}_y, \hat{y}] = \hat{x}(-\iu\hbar\hat{I}) = \boxed{-\iu\hbar\hat{x}}.
        \label{eq:lz_y_comm}
    \end{equation}

    \textit{Paso 2: Expansión de Hadamard:}
    Como $\hat{L}_z$ es hermítico, el operador conjugante es $\hat{R}_z\adj(\theta) = \exp\left(\frac{\iu}{\hbar}\theta\hat{L}_z\right)$. Aplicamos el Lema de Hadamard definiendo el operador $\hat{A} \equiv \frac{\iu\theta}{\hbar}\hat{L}_z$ y $\hat{B} \equiv \hat{x}$:
    \begin{equation}
        \hat{x}' = \eu^{\hat{A}}\hat{x}\eu^{-\hat{A}} = \hat{x} + [\hat{A}, \hat{x}] + \frac{1}{2!}[\hat{A}, [\hat{A}, \hat{x}]] + \frac{1}{3!}[\hat{A}, [\hat{A}, [\hat{A}, \hat{x}]]] + \dots
    \end{equation}
    Evaluemos término a término los conmutadores anidados:
    \begin{itemize}
        \item Primer orden ($n = 1$):
        \begin{equation}
            [\hat{A}, \hat{x}] = \left[\frac{\iu\theta}{\hbar}\hat{L}_z, \, \hat{x}\right] = \frac{\iu\theta}{\hbar}[\hat{L}_z, \hat{x}] = \frac{\iu\theta}{\hbar}(\iu\hbar\hat{y}) = \iu^2 \theta \hat{y} = -\theta\hat{y}.
        \end{equation}
        \item Segundo orden ($n = 2$):
        \begin{equation}
            [\hat{A}, [\hat{A}, \hat{x}]] = [\hat{A}, -\theta\hat{y}] = -\theta\left[\frac{\iu\theta}{\hbar}\hat{L}_z, \, \hat{y}\right] = -\frac{\iu\theta^2}{\hbar}(-\iu\hbar\hat{x}) = \iu^2 \theta^2 \hat{x} = -\theta^2 \hat{x}.
        \end{equation}
        \item Tercer orden ($n = 3$):
        \begin{equation}
            [\hat{A}, [\hat{A}, [\hat{A}, \hat{x}]]] = [\hat{A}, -\theta^2\hat{x}] = -\theta^2 [\hat{A}, \hat{x}] = -\theta^2(-\theta\hat{y}) = +\theta^3\hat{y}.
        \end{equation}
        \item Cuarto orden ($n = 4$):
        \begin{equation}
            [\hat{A}, [\hat{A}, [\hat{A}, [\hat{A}, \hat{x}]]]] = [\hat{A}, \theta^3\hat{y}] = \theta^3(+\theta\hat{x}) = +\theta^4\hat{x}.
        \end{equation}
    \end{itemize}
    Sustituyendo estos resultados en la serie infinita y agrupando por las componentes de los operadores $\hat{x}$ y $\hat{y}$:
    \begin{align}
        \hat{x}' &= \hat{x}\left(1 - \frac{\theta^2}{2!} + \frac{\theta^4}{4!} - \dots\right) - \hat{y}\left(\theta - \frac{\theta^3}{3!} + \frac{\theta^5}{5!} - \dots\right).
    \end{align}
    Reconocemos de inmediato las expansiones de Taylor clásicas de las funciones coseno y seno:
    \begin{equation}
        \boxed{\hat{x}' = \hat{x}\cos\theta - \hat{y}\sin\theta}.
    \end{equation}
    De manera idéntica se demuestra que $\hat{y}' = \hat{x}\sin\theta + \hat{y}\cos\theta$. ¡El Lema de Hadamard genera la rotación geométrica exacta en el plano euclídeo a partir del álgebra pura de conmutadores!
\end{ejemplobox}

\begin{figure}[H]
\centering
\begin{subfigure}[t]{0.45\textwidth}
\centering
\begin{tikzpicture}[scale=2.2]
  \draw[->,gray] (-0.3,0) -- (1.35,0) node[right] {$\hat x$};
  \draw[->,gray] (0,-0.9) -- (0,1.2) node[above] {$\hat y$};
  \draw[dashed,gray!60] (0,0) circle (1);
  \draw[-{Stealth},very thick,uniblue] (0,0) -- (1,0);
  \draw[-{Stealth},very thick,unired] (0,0) -- ({cos(40)},{-sin(40)})
        node[below right] {$\hat x'$};
  \draw[unigreen,thick] (0.45,0) arc (0:-40:0.45);
  \node[unigreen] at (0.62,-0.13) {$\theta$};
  \draw[dotted,unired] ({cos(40)},{-sin(40)}) -- ({cos(40)},0);
  \draw[dotted,unired] ({cos(40)},{-sin(40)}) -- (0,{-sin(40)});
  \node[font=\scriptsize,unired,above] at ({cos(40)},0) {$\cos\theta$};
  \node[font=\scriptsize,unired,left] at (0,{-sin(40)}) {$-\sin\theta$};
\end{tikzpicture}
\caption{$\hat x'=\hat x\cos\theta-\hat y\sin\theta$ (componentes sobre $\hat x,\hat y$).}
\end{subfigure}\hfill
\begin{subfigure}[t]{0.52\textwidth}
\centering
\begin{tikzpicture}
\begin{axis}[width=\linewidth,height=6.2cm,xlabel={$\theta$},
  xmin=0,xmax=7,ymin=-2.2,ymax=2.2,
  legend style={font=\scriptsize,at={(0.5,-0.28)},anchor=north,legend columns=3},
  grid=both,minor tick num=1,
  xtick={0,1.5708,3.1416,4.7124,6.2832},
  xticklabels={$0$,$\pi/2$,$\pi$,$3\pi/2$,$2\pi$}]
  \addplot[black,very thick,domain=0:7,samples=200] {cos(deg(x))};
  \addlegendentry{$\cos\theta$}
  \addplot[uniblue,dashed,thick,domain=0:7,samples=100] {1};
  \addlegendentry{orden 0}
  \addplot[unired,thick,domain=0:7,samples=100] {1-x^2/2};
  \addlegendentry{hasta orden 2}
  \addplot[unigreen,thick,domain=0:7,samples=100] {1-x^2/2+x^4/24};
  \addlegendentry{hasta orden 4}
  \addplot[orange,thick,domain=0:7,samples=100] {1-x^2/2+x^4/24-x^6/720};
  \addlegendentry{hasta orden 6}
\end{axis}
\end{tikzpicture}
\caption{Sumas parciales de la serie de Hadamard (coeficiente de $\hat x$).}
\end{subfigure}
\caption{Rotaci\'on finita generada por $\hat L_z$: cada orden de la serie de Hadamard a\~nade un t\'ermino de la serie de Taylor de $\cos\theta$ y $\sin\theta$.}
\label{fig:rot_hadamard}
\end{figure}

\subsubsection[Problema 2: Operador vectorial y tensores]{Problema Resuelto 2: El Operador Vectorial Cuántico y Tensores Cartesianos}
\begin{ejemplobox}[Problema Resuelto 2: El Operador Vectorial Cuántico]
    \textbf{Enunciado:} En mecánica cuántica, se define formalmente como \textbf{operador vectorial} a cualquier triplete de observables $\hat{\vec{V}} = (\hat{V}_1, \hat{V}_2, \hat{V}_3)$ que bajo una rotación espacial se transforme como un vector geométrico, lo cual equivale a la relación de conmutación:
    \begin{equation}
        [\hat{L}_i, \hat{V}_j] = \iu\hbar\,\varepsilon_{ijk}\hat{V}_k.
    \end{equation}
    Demuestre que el vector de posición $\hat{\vec{r}} = (\hat{x}_1, \hat{x}_2, \hat{x}_3)$ y el vector de momento lineal $\hat{\vec{p}} = (\hat{p}_1, \hat{p}_2, \hat{p}_3)$ son operadores vectoriales genuinos empleando cálculo tensorial.

    \textbf{Solución Paso a Paso:}
    
    Escribimos la componente $\hat{L}_i$ en notación tensorial: $\hat{L}_i = \varepsilon_{iab}\hat{x}_a\hat{p}_b$. Evaluamos el conmutador con la coordenada $\hat{x}_j$:
    \begin{align}
        [\hat{L}_i, \hat{x}_j] &= [\varepsilon_{iab}\hat{x}_a\hat{p}_b, \, \hat{x}_j] = \varepsilon_{iab}[\hat{x}_a\hat{p}_b, \, \hat{x}_j].
    \end{align}
    Aplicando la regla del conmutador de un producto $[\hat{A}\hat{B}, \hat{C}] = \hat{A}[\hat{B}, \hat{C}] + [\hat{A}, \hat{C}]\hat{B}$:
    \begin{equation}
        [\hat{x}_a\hat{p}_b, \, \hat{x}_j] = \hat{x}_a[\hat{p}_b, \hat{x}_j] + [\hat{x}_a, \hat{x}_j]\hat{p}_b.
    \end{equation}
    Dado que $[\hat{x}_a, \hat{x}_j] = 0$ y $[\hat{p}_b, \hat{x}_j] = -\iu\hbar\delta_{bj}$:
    \begin{equation}
        [\hat{x}_a\hat{p}_b, \, \hat{x}_j] = \hat{x}_a(-\iu\hbar\delta_{bj}) = -\iu\hbar\delta_{bj}\hat{x}_a.
    \end{equation}
    Sustituyendo esto dentro de la suma tensorial:
    \begin{equation}
        [\hat{L}_i, \hat{x}_j] = \varepsilon_{iab}(-\iu\hbar\delta_{bj}\hat{x}_a) = -\iu\hbar\,\varepsilon_{iaj}\hat{x}_a.
    \end{equation}
    Intercambiamos los dos últimos índices del tensor de Levi-Civita para invertir el signo: $\varepsilon_{iaj} = -\varepsilon_{ija}$:
    \begin{equation}
        [\hat{L}_i, \hat{x}_j] = \iu\hbar\,\varepsilon_{ija}\hat{x}_a = \boxed{\iu\hbar\,\varepsilon_{ijk}\hat{x}_k}.
    \end{equation}
    Para el momento $\hat{p}_j$, evaluamos de la misma manera:
    \begin{equation}
        [\hat{L}_i, \hat{p}_j] = \varepsilon_{iab}[\hat{x}_a\hat{p}_b, \, \hat{p}_j] = \varepsilon_{iab}[\hat{x}_a, \hat{p}_j]\hat{p}_b = \varepsilon_{iab}(\iu\hbar\delta_{aj})\hat{p}_b = \iu\hbar\,\varepsilon_{ijb}\hat{p}_b = \boxed{\iu\hbar\,\varepsilon_{ijk}\hat{p}_k}.
    \end{equation}
    Ambos operadores satisfacen la definición exacta de operador vectorial cuántico.
\end{ejemplobox}

\subsubsection[Problema 3: Incertidumbre de momento angular en $\ket{\ell, m}$]{Problema Resuelto 3: Incertidumbre de Momento Angular en $\ket{\ell, m}$ y Operadores Escalera}
\begin{ejemplobox}[Problema Resuelto 3: Incertidumbre de Momento Angular en el Autoestado $\ket{\ell, m}$]
    \textbf{Enunciado:} Considere una partícula en un autoestado común $\ket{\ell, m}$ de $\hat{L}^2$ y $\hat{L}_z$.
    \begin{enumerate}
        \item Calcule los valores esperados $\ev{\hat{L}_x}$ y $\ev{\hat{L}_y}$.
        \item Calcule $\ev{\hat{L}_x^2}$ y $\ev{\hat{L}_y^2}$ utilizando los operadores escalera $\hat{L}_\pm = \hat{L}_x \pm \iu\hat{L}_y$.
        \item Verifique que el producto de incertidumbres $\Delta L_x \Delta L_y$ cumple estrictamente la relación de Heisenberg.
    \end{enumerate}

    \textbf{Solución Paso a Paso:}

    \textit{Paso 1: Operadores Escalera y Valores Esperados:}
    Los operadores escalera se definen como $\hat{L}_+ = \hat{L}_x + \iu\hat{L}_y$ y $\hat{L}_- = \hat{L}_x - \iu\hat{L}_y$. Despejando las componentes cartesianas:
    \begin{equation}
        \hat{L}_x = \frac{1}{2}(\hat{L}_+ + \hat{L}_-), \qquad \hat{L}_y = \frac{1}{2\iu}(\hat{L}_+ - \hat{L}_-).
    \end{equation}
    La acción de los operadores escalera sobre la base ortonormal es:
    \begin{equation}
        \hat{L}_\pm \ket{\ell, m} = \hbar\sqrt{\ell(\ell+1) - m(m\pm 1)}\ket{\ell, m\pm 1}.
    \end{equation}
    Al tomar el producto interno con $\bra{\ell, m}$, como $\braket{\ell, m}{\ell, m\pm 1} = 0$ por ortogonalidad:
    \begin{equation}
        \ev{\hat{L}_+} = \matrixel{\ell, m}{\hat{L}_+}{\ell, m} = 0, \qquad \ev{\hat{L}_-} = \matrixel{\ell, m}{\hat{L}_-}{\ell, m} = 0.
    \end{equation}
    Por lo tanto:
    \begin{equation}
        \ev{\hat{L}_x} = 0, \qquad \ev{\hat{L}_y} = 0.
    \end{equation}

    \textit{Paso 2: Varianza Cuántica y Demostración Algebraica de $\ev{\hat{L}_x^2} = \ev{\hat{L}_y^2}$:}
    Para evaluar rigurosamente $\ev{\hat{L}_x^2}$ y $\ev{\hat{L}_y^2}$ sin recurrir a argumentos heurísticos de simetría, expresemos los cuadrados directamente en términos de los operadores escalera:
    \begin{align}
        \hat{L}_x^2 &= \frac{1}{4}\left(\hat{L}_+ + \hat{L}_-\right)^2 = \frac{1}{4}\left( \hat{L}_+^2 + \hat{L}_-^2 + \hat{L}_+\hat{L}_- + \hat{L}_-\hat{L}_+ \right), \\
        \hat{L}_y^2 &= -\frac{1}{4}\left(\hat{L}_+ - \hat{L}_-\right)^2 = -\frac{1}{4}\left( \hat{L}_+^2 + \hat{L}_-^2 - \hat{L}_+\hat{L}_- - \hat{L}_-\hat{L}_+ \right).
    \end{align}
    Al tomar el valor esperado en el estado $\ket{\ell, m}$, los términos $\hat{L}_\pm^2$ desplazan el número cuántico magnético en dos unidades ($\hat{L}_\pm^2\ket{\ell, m} \propto \ket{\ell, m \pm 2}$). Como la base es ortogonal ($\braket{\ell, m}{\ell, m \pm 2} = 0$), sus valores esperados se anulan idénticamente:
    \begin{equation}
        \ev{\hat{L}_+^2} = \matrixel{\ell, m}{\hat{L}_+^2}{\ell, m} = 0, \qquad \ev{\hat{L}_-^2} = \matrixel{\ell, m}{\hat{L}_-^2}{\ell, m} = 0.
    \end{equation}
    En consecuencia, ambos valores esperados se reducen de forma exacta a la misma expresión simétrica:
    \begin{equation}
        \ev{\hat{L}_x^2} = \ev{\hat{L}_y^2} = \frac{1}{4}\ev{\hat{L}_+\hat{L}_- + \hat{L}_-\hat{L}_+}.
    \end{equation}
    Por otro lado, la suma de las tres componentes al cuadrado es el operador Casimir $\hat{L}^2$:
    \begin{equation}
        \hat{L}^2 = \hat{L}_x^2 + \hat{L}_y^2 + \hat{L}_z^2 \implies \ev{\hat{L}_x^2} + \ev{\hat{L}_y^2} = \ev{\hat{L}^2} - \ev{\hat{L}_z^2} = \hbar^2[\ell(\ell+1) - m^2].
    \end{equation}
    Como ambos valores esperados son iguales, se deduce algebraicamente:
    \begin{equation}
        \ev{\hat{L}_x^2} = \ev{\hat{L}_y^2} = \frac{\hbar^2}{2}[\ell(\ell+1) - m^2].
    \end{equation}
    Dado que las medias son nulas ($\ev{\hat{L}_x} = \ev{\hat{L}_y} = 0$), las desviaciones estándar son:
    \begin{equation}
        \Delta L_x = \Delta L_y = \hbar\sqrt{\frac{\ell(\ell+1) - m^2}{2}}.
    \end{equation}

    \textit{Paso 3: Verificación de la Relación de Incertidumbre:}
    El producto de incertidumbres resulta ser:
    \begin{equation}
        \Delta L_x \Delta L_y = \frac{\hbar^2}{2}[\ell(\ell+1) - m^2].
    \end{equation}
    Por otro lado, la cota inferior dictada por la relación de Heisenberg para $[\hat{L}_x, \hat{L}_y] = \iu\hbar\hat{L}_z$ es:
    \begin{equation}
        \frac{1}{2}|\ev{[\hat{L}_x, \hat{L}_y]}| = \frac{\hbar}{2}|\ev{\hat{L}_z}| = \frac{\hbar^2}{2}|m|.
    \end{equation}
    Comprobemos que este valor estadístico satisface rigurosamente el principio de incertidumbre de Heisenberg generalizado:
    \begin{equation}
        \frac{\hbar^2}{2}[\ell(\ell+1) - m^2] \ge \frac{\hbar^2}{2}|m|.
    \end{equation}
    Esto equivale algebraicamente a:
    \begin{equation}
        \ell(\ell+1) - m^2 \ge |m| \iff \ell(\ell+1) \ge m^2 + |m| = |m|(|m|+1).
    \end{equation}
    Dado que por la teoría general de momentos angulares $|m| \le \ell$, la función $f(x) = x(x+1)$ es monótonamente creciente para $x \ge 0$. Por consiguiente, $\ell(\ell+1) \ge |m|(|m|+1)$ se satisface para todo $m \in \{-\ell, \dots, \ell\}$. La igualdad se satura si y solo si $|m| = \ell$ (estados de máximo alineamiento angular).
\end{ejemplobox}

\begin{figure}[H]
    \centering
    % Fig. 4: producto de incertidumbres vs cota, \ell = 3 (unidades hbar^2)
\begin{tikzpicture}
\begin{axis}[width=0.8\linewidth, height=6cm, domain=-3:3, samples=100,
  xlabel={$m$}, ylabel={$\Delta L_x\,\Delta L_y\;/\;\hbar^2$},
  axis lines=left, xtick={-3,-2,-1,0,1,2,3}, ymin=0, ymax=9.5,
  legend style={font=\scriptsize, draw=none, fill=white, at={(0.5,0.97)}, anchor=north},
  every axis plot/.append style={thick}]
  \addplot[uni-blue]{0.5*(12-x^2)};
  \addplot[alert-border]{0.5*abs(x)};
  \addplot[only marks, mark=*, uni-blue, mark size=2pt, samples at={-3,-2,...,3}]{0.5*(12-x^2)};
  \addplot[only marks, mark=square*, example-border, mark size=2.5pt, samples at={-3,3}]{0.5*abs(x)};
  \legend{$\tfrac12[\ell(\ell+1)-m^2]$, cota $\tfrac12|m|$, valores permitidos, saturaci\'on ($|m|=\ell$)}
\end{axis}
\end{tikzpicture}

    \caption{Producto $\Delta L_x\,\Delta L_y$ en $\ket{\ell,m}$ frente a la cota de Heisenberg, para $\ell=3$. La igualdad solo se alcanza en $|m|=\ell$.}
    \label{fig:incert_Lxy}
\end{figure}

\subsubsection[Problema 4: Efecto Zeeman y precesión de Larmor]{Problema Resuelto 4: Efecto Zeeman Orbital, Ecuaciones de Heisenberg y Precesión de Larmor}
\begin{ejemplobox}[Problema Resuelto 4: Efecto Zeeman Orbital y Precesión de Larmor Cuántica]
    \textbf{Enunciado:} Un electrón de masa $m_e$ y carga $-e$ en un átomo se encuentra inmerso en un campo magnético estático y homogéneo $\vec{B} = B_0\hat{z}$. El momento dipolar magnético orbital está acoplado al momento angular orbital mediante $\hat{\vec{\mu}}_L = -\frac{e}{2m_e}\hat{\vec{L}}$. El Hamiltoniano total del sistema (efecto Zeeman normal) viene dado por:
    \begin{equation}
        \hat{H} = \hat{H}_0 - \hat{\vec{\mu}}_L \cdot \vec{B} = \hat{H}_0 + \omega_L \hat{L}_z,
    \end{equation}
    donde $\omega_L \equiv \frac{e B_0}{2m_e}$ es la frecuencia clásica de Larmor y $[\hat{H}_0, \hat{\vec{L}}] = 0$.
    \begin{enumerate}
        \item Calcule los conmutadores fundamentales $[\hat{H}, \hat{L}_x]$, $[\hat{H}, \hat{L}_y]$ y $[\hat{H}, \hat{L}_z]$.
        \item En la representación de Heisenberg, la evolución temporal de cualquier observable satisface $\frac{\dif\hat{A}_H}{\dif t} = \frac{1}{\iu\hbar}[\hat{A}_H, \hat{H}]$. Deduzca el sistema de ecuaciones diferenciales acopladas para $\frac{\dif\hat{L}_x}{\dif t}$, $\frac{\dif\hat{L}_y}{\dif t}$ y $\frac{\dif\hat{L}_z}{\dif t}$.
        \item Desacople el sistema definiendo los operadores escalera $\hat{L}_\pm = \hat{L}_x \pm \iu\hat{L}_y$ y obtenga las soluciones analíticas para $\hat{L}_x(t)$ y $\hat{L}_y(t)$ en función de sus valores iniciales en $t = 0$.
        \item Demuestre el Teorema de Ehrenfest para este sistema y describa la trayectoria cuántica de los valores esperados $\ev{\hat{\vec{L}}(t)}$.
    \end{enumerate}

    \textbf{Solución Paso a Paso:}

    \textit{Paso 1: Conmutadores con el Hamiltoniano:}
    Puesto que $[\hat{H}_0, \hat{\vec{L}}] = 0$, los conmutadores se reducen al acoplamiento con $\hat{L}_z$:
    \begin{align}
        [\hat{H}, \hat{L}_x] &= [\hat{H}_0 + \omega_L\hat{L}_z, \, \hat{L}_x] = \omega_L [\hat{L}_z, \hat{L}_x] = \omega_L(\iu\hbar\hat{L}_y) = \iu\hbar\omega_L\hat{L}_y, \\
        [\hat{H}, \hat{L}_y] &= [\hat{H}_0 + \omega_L\hat{L}_z, \, \hat{L}_y] = \omega_L [\hat{L}_z, \hat{L}_y] = \omega_L(-\iu\hbar\hat{L}_x) = -\iu\hbar\omega_L\hat{L}_x, \\
        [\hat{H}, \hat{L}_z] &= \omega_L [\hat{L}_z, \hat{L}_z] = 0.
    \end{align}

    \textit{Paso 2: Ecuaciones de Movimiento de Heisenberg:}
    Sustituyendo los conmutadores en la ecuación de Heisenberg (notando que $[\hat{L}_i, \hat{H}] = -[\hat{H}, \hat{L}_i]$):
    \begin{align}
        \frac{\dif\hat{L}_x}{\dif t} &= \frac{1}{\iu\hbar}[\hat{L}_x, \hat{H}] = -\frac{1}{\iu\hbar}[\hat{H}, \hat{L}_x] = -\frac{1}{\iu\hbar}(\iu\hbar\omega_L\hat{L}_y) = -\omega_L\hat{L}_y, \\
        \frac{\dif\hat{L}_y}{\dif t} &= \frac{1}{\iu\hbar}[\hat{L}_y, \hat{H}] = -\frac{1}{\iu\hbar}[\hat{H}, \hat{L}_y] = -\frac{1}{\iu\hbar}(-\iu\hbar\omega_L\hat{L}_x) = +\omega_L\hat{L}_x, \\
        \frac{\dif\hat{L}_z}{\dif t} &= \frac{1}{\iu\hbar}[\hat{L}_z, \hat{H}] = 0 \implies \hat{L}_z(t) = \hat{L}_z(0) = \text{Constante de movimiento}.
    \end{align}

    \textit{Paso 3: Desacoplamiento mediante Operadores Escalera:}
    Combinando linealmente las dos primeras ecuaciones:
    \begin{equation}
        \frac{\dif\hat{L}_\pm}{\dif t} = \frac{\dif\hat{L}_x}{\dif t} \pm \iu\frac{\dif\hat{L}_y}{\dif t} = -\omega_L\hat{L}_y \pm \iu\omega_L\hat{L}_x = \pm\iu\omega_L(\hat{L}_x \pm \iu\hat{L}_y) = \pm\iu\omega_L\hat{L}_\pm.
    \end{equation}
    ¡Las ecuaciones para $\hat{L}_+$ y $\hat{L}_-$ están completamente desacopladas! Integrando directamente respecto al tiempo:
    \begin{equation}
        \hat{L}_\pm(t) = \hat{L}_\pm(0)\eu^{\pm\iu\omega_L t}.
    \end{equation}
    Recuperamos las componentes cartesianas usando $\hat{L}_x = \frac{1}{2}(\hat{L}_+ + \hat{L}_-)$ y $\hat{L}_y = \frac{1}{2\iu}(\hat{L}_+ - \hat{L}_-)$:
    \begin{align}
        \hat{L}_x(t) &= \frac{1}{2}\left( \hat{L}_+(0)\eu^{\iu\omega_L t} + \hat{L}_-(0)\eu^{-\iu\omega_L t} \right) \nonumber\\
        &= \frac{1}{2}\left( (\hat{L}_x(0) + \iu\hat{L}_y(0))\eu^{\iu\omega_L t} + (\hat{L}_x(0) - \iu\hat{L}_y(0))\eu^{-\iu\omega_L t} \right) \nonumber\\
        &= \hat{L}_x(0)\left(\frac{\eu^{\iu\omega_L t} + \eu^{-\iu\omega_L t}}{2}\right) - \hat{L}_y(0)\left(\frac{\eu^{\iu\omega_L t} - \eu^{-\iu\omega_L t}}{2\iu}\right) \nonumber\\
        &= \boxed{\hat{L}_x(0)\cos(\omega_L t) - \hat{L}_y(0)\sin(\omega_L t)}.
    \end{align}
    Análogamente, para la componente transversal $y$:
    \begin{equation}
        \boxed{\hat{L}_y(t) = \hat{L}_x(0)\sin(\omega_L t) + \hat{L}_y(0)\cos(\omega_L t)}.
    \end{equation}

    \textit{Paso 4: Teorema de Ehrenfest y Precesión de Larmor:}
    Tomando el valor esperado en un estado cuántico arbitrario $\ket{\psi}$:
    \begin{equation}
        \frac{\dif}{\dif t}\ev{\hat{\vec{L}}} = \vec{\omega}_L \times \ev{\hat{\vec{L}}}, \quad \text{con } \vec{\omega}_L = \omega_L\hat{z}.
    \end{equation}
    El vector valor esperado $\ev{\hat{\vec{L}}(t)}$ describe una rotación circular uniforme en el plano $xy$ con velocidad angular constante $\omega_L$, manteniendo constante su ángulo de inclinación respecto al eje $z$. ¡La mecánica cuántica reproduce con total fidelidad la célebre precesión de Larmor del giróscopo magnético clásico!
\end{ejemplobox}

\subsubsection[Problema 5: Oscilador 2D y modos circulares]{Problema Resuelto 5: Oscilador Armónico 2D Isótropo y Modos Circulares \texorpdfstring{$\ket{R}$}{R} y \texorpdfstring{$\ket{L}$}{L}}
\begin{ejemplobox}[Problema Resuelto 5: Oscilador 2D Isótropo y Modos Circulares Cuánticos]
    \textbf{Enunciado:} Considere una partícula de masa $M$ en un potencial armónico bidimensional isótropo:
    \begin{equation}
        \hat{H} = \frac{\hat{p}_x^2 + \hat{p}_y^2}{2M} + \frac{1}{2}M\omega^2(\hat{x}^2 + \hat{y}^2) = \hbar\omega\left( \hat{a}_x\adj\hat{a}_x + \hat{a}_y\adj\hat{a}_y + 1 \right),
    \end{equation}
    y el operador momento angular orbital perpendicular al plano de oscilación:
    \begin{equation}
        \hat{L}_z = \hat{x}\hat{p}_y - \hat{y}\hat{p}_x = -\iu\hbar\left(\hat{a}_x\adj\hat{a}_y - \hat{a}_y\adj\hat{a}_x\right).
    \end{equation}
    \begin{enumerate}
        \item Demuestre que $[\hat{H}, \hat{L}_z] = 0$. ¿Qué garantiza este resultado sobre los autoestados de energía?
        \item Para el primer nivel de energía excitado $E_1 = 2\hbar\omega$, la base cartesiana desacoplada $\{\ket{n_x, n_y}\}$ es bidimensional: $\{\ket{1, 0} = \hat{a}_x\adj\ket{0, 0}, \, \ket{0, 1} = \hat{a}_y\adj\ket{0, 0}\}$. Construya la matriz de $2\times 2$ representante de $\hat{L}_z$ restringida a este subespacio degenerado $E_1$.
        \item Diagonalice la matriz de $\hat{L}_z$ en $E_1$ y halle los autovalores y autoestados simultáneos comunes a $\hat{H}$ y $\hat{L}_z$.
        \item Demuestre que estos autoestados corresponden a modos circulares cuánticos dextrógiro $\ket{R}$ y levógiro $\ket{L}$ con autovalores $L_z = \pm\hbar$.
    \end{enumerate}

    \textbf{Solución Paso a Paso:}

    \textit{Paso 1: Conmutación $[\hat{H}, \hat{L}_z] = 0$:}
    Escribimos $\hat{H} = \hbar\omega(\hat{N} + 1)$, donde $\hat{N} = \hat{a}_x\adj\hat{a}_x + \hat{a}_y\adj\hat{a}_y$. Evaluamos el conmutador con $\hat{L}_z$:
    \begin{align}
        [\hat{N}, \hat{a}_x\adj\hat{a}_y] &= [\hat{a}_x\adj\hat{a}_x, \hat{a}_x\adj\hat{a}_y] + [\hat{a}_y\adj\hat{a}_y, \hat{a}_x\adj\hat{a}_y] \nonumber\\
        &= \hat{a}_x\adj[\hat{a}_x, \hat{a}_x\adj]\hat{a}_y + \hat{a}_x\adj[\hat{a}_y\adj\hat{a}_y, \hat{a}_y] \nonumber\\
        &= \hat{a}_x\adj(1)\hat{a}_y + \hat{a}_x\adj(-\hat{a}_y) = \hat{a}_x\adj\hat{a}_y - \hat{a}_x\adj\hat{a}_y = 0.
    \end{align}
    Análogamente, $[\hat{N}, \hat{a}_y\adj\hat{a}_x] = 0$. Por linealidad:
    \begin{equation}
        [\hat{H}, \hat{L}_z] = -\iu\hbar^2\omega [\hat{N}, \, \hat{a}_x\adj\hat{a}_y - \hat{a}_y\adj\hat{a}_x] = 0.
    \end{equation}
    Por el Teorema 2.3, $\hat{H}$ y $\hat{L}_z$ son compatibles y admiten una base ortonormal común.

    \textit{Paso 2: Matriz Representante en el Subespacio Degenerado $E_1$:}
    La base ordenada de $E_1$ es $\{\ket{1, 0}, \ket{0, 1}\}$. Evaluemos la acción de $\hat{L}_z = -\iu\hbar(\hat{a}_x\adj\hat{a}_y - \hat{a}_y\adj\hat{a}_x)$:
    \begin{align}
        \hat{L}_z\ket{1, 0} &= -\iu\hbar\left( \hat{a}_x\adj\hat{a}_y\ket{1, 0} - \hat{a}_y\adj\hat{a}_x\ket{1, 0} \right) = -\iu\hbar(0 - \hat{a}_y\adj\ket{0, 0}) = +\iu\hbar\ket{0, 1}, \\
        \hat{L}_z\ket{0, 1} &= -\iu\hbar\left( \hat{a}_x\adj\hat{a}_y\ket{0, 1} - \hat{a}_y\adj\hat{a}_x\ket{0, 1} \right) = -\iu\hbar(\hat{a}_x\adj\ket{0, 0} - 0) = -\iu\hbar\ket{1, 0}.
    \end{align}
    Calculamos los elementos de matriz en la base $\{\ket{1, 0}, \ket{0, 1}\}$:
    \begin{equation}
        \hat{L}_z|_{E_1} = \begin{pmatrix}
            \matrixel{1, 0}{\hat{L}_z}{1, 0} & \matrixel{1, 0}{\hat{L}_z}{0, 1} \\
            \matrixel{0, 1}{\hat{L}_z}{1, 0} & \matrixel{0, 1}{\hat{L}_z}{0, 1}
        \end{pmatrix} = \hbar \begin{pmatrix} 0 & -\iu \\ \iu & 0 \end{pmatrix} = \hbar\sigma_y.
    \end{equation}
    ¡La matriz representante en este subespacio es exactamente la matriz de Pauli $\sigma_y$ multiplicada por $\hbar$!

    \textit{Paso 3 y 4: Diagonalización y Modos Circulares $\ket{R}$ y $\ket{L}$:}
    El polinomio característico es:
    \begin{equation}
        \det(\hat{L}_z|_{E_1} - \lambda I) = \begin{vmatrix} -\lambda & -\iu\hbar \\ \iu\hbar & -\lambda \end{vmatrix} = \lambda^2 - \hbar^2 = 0 \implies \lambda = \pm\hbar.
    \end{equation}
    Determinemos los autovectores normalizados para cada autovalor:
    \begin{itemize}
        \item Para $\lambda_1 = +\hbar$ (modo circular dextrógiro $\ket{R}$):
        \begin{equation}
            \begin{pmatrix} -\hbar & -\iu\hbar \\ \iu\hbar & -\hbar \end{pmatrix}\begin{pmatrix} c_1 \\ c_2 \end{pmatrix} = 0 \implies c_1 = -\iu c_2 \implies \boxed{\ket{R} = \frac{1}{\sqrt{2}}\left( \ket{1, 0} + \iu\ket{0, 1} \right)}.
        \end{equation}
        \item Para $\lambda_2 = -\hbar$ (modo circular levógiro $\ket{L}$):
        \begin{equation}
            \begin{pmatrix} \hbar & -\iu\hbar \\ \iu\hbar & \hbar \end{pmatrix}\begin{pmatrix} c_1 \\ c_2 \end{pmatrix} = 0 \implies c_1 = \iu c_2 \implies \boxed{\ket{L} = \frac{1}{\sqrt{2}}\left( \ket{1, 0} - \iu\ket{0, 1} \right)}.
        \end{equation}
    \end{itemize}
    Notemos la física detrás del resultado: el estado $\ket{1, 0}$ representa una oscilación lineal pura a lo largo del eje $x$, mientras que $\ket{0, 1}$ oscila a lo largo de $y$. La combinación $\ket{R} \propto \ket{1, 0} + \iu\ket{0, 1}$ introduce un desfase cuántico de $\pi/2$ ($\eu^{\iu\pi/2} = \iu$), exactamente idéntico a cómo dos ondas clásicas desfasadas en cuadratura generan polarización circular dextrógira y levógira. ¡El Teorema 2.3 (diagonalización simultánea) en bloques degenerados resuelve el problema con elegancia absoluta!
\end{ejemplobox}

\subsection{Ejercicios Propuestos para Autoestudio (Nivel Avanzado / Zettili)}

\subsubsection[Ejercicio 1: Conmutador de la exponencial]{Ejercicio Propuesto 1: Conmutador de la Exponencial de un Operador}
\begin{ejerciciobox}[Ejercicio Propuesto 1: Conmutador de la Exponencial de un Operador]
    \textbf{Nivel: Intermedio}\\
    Sean $\hat{A}$ y $\hat{B}$ dos operadores que satisfacen la condición $[[\hat{A}, \hat{B}], \hat{A}] = 0$.
    \begin{enumerate}
        \item Demuestre que $[\hat{A}^n, \hat{B}] = n\hat{A}^{n-1}[\hat{A}, \hat{B}]$ para todo entero positivo $n \ge 1$.
        \item Utilice la expansión en serie de potencias de la exponencial para demostrar formalmente que:
        \begin{equation}
            [\eu^{\hat{A}}, \hat{B}] = [\hat{A}, \hat{B}]\eu^{\hat{A}}.
        \end{equation}
    \end{enumerate}
    \textit{Pista:} Utilice inducción matemática para la parte (1) y derive término a término en la parte (2).
\end{ejerciciobox}

\subsubsection[Ejercicio 2: Generador de dilataciones $\hat{D}$]{Ejercicio Propuesto 2: Invariancia Rotacional del Generador de Dilataciones $\hat{D}$}
\begin{ejerciciobox}[Ejercicio Propuesto 2: Invariancia Rotacional del Generador de Dilataciones $\hat{D}$]
    \textbf{Nivel: Avanzado}\\
    El generador hermítico de las transformaciones de escala (dilataciones espaciales) en tres dimensiones se define mediante la combinación simetrizada de posición y momento:
    \begin{equation}
        \hat{D} \equiv \frac{1}{2}\left(\hat{\vec{r}}\cdot\hat{\vec{p}} + \hat{\vec{p}}\cdot\hat{\vec{r}}\right) = \hat{\vec{r}}\cdot\hat{\vec{p}} - \frac{3}{2}\iu\hbar\hat{I}.
    \end{equation}
    \begin{enumerate}
        \item Verifique que $\hat{D}$ es estrictamente autoadjunto ($\hat{D}\adj = \hat{D}$) y demuestre que conmuta con cada componente del momento angular: $[\hat{L}_i, \hat{D}] = 0$ para todo $i \in \{1, 2, 3\}$.
        \item Demuestre que el operador Casimir $\hat{L}^2$ conmuta con $\hat{D}$:
        \begin{equation}
            [\hat{L}^2, \, \hat{D}] = 0.
        \end{equation}
        ¿Cuál es el significado físico de este resultado en términos de simetrías continuas y transformaciones de escala?
    \end{enumerate}
    \textit{Intuición Didáctica y Pista:} En la física clásica, la cantidad escalar que mide dilataciones de escala es $\vec{r}\cdot\vec{p}$. No obstante, en mecánica cuántica, como $\hat{x}_k$ y $\hat{p}_k$ no conmutan, el producto simple no es hermítico: $(\hat{\vec{r}}\cdot\hat{\vec{p}})\adj = \hat{\vec{r}}\cdot\hat{\vec{p}} - 3\iu\hbar\hat{I}$. La \textbf{regla de oro} de la cuantización canónica para construir observables físicos a partir de productos clásicos es promediar simétricamente ambos órdenes: $\frac{1}{2}(\hat{A}\hat{B} + \hat{B}\hat{A})$ (conocido como ordenamiento de Weyl). Como la constante $-\frac{3}{2}\iu\hbar\hat{I}$ conmuta con todo operador, concéntrese en evaluar $[\hat{L}_i, \hat{x}_k\hat{p}_k] = [\hat{L}_i, \hat{x}_k]\hat{p}_k + \hat{x}_k[\hat{L}_i, \hat{p}_k]$ y contraiga el tensor simétrico resultante con el Levi-Civita $\varepsilon_{ikm}$.
\end{ejerciciobox}

\subsubsection[Ejercicio 3: Robertson--Schrödinger y chirp]{Ejercicio Propuesto 3: Relación de Robertson--Schrödinger en Paquetes con Chirp}
\begin{ejerciciobox}[Ejercicio Propuesto 3: Relación de Robertson--Schrödinger en Estados Coherentes Comprimidos]
    \textbf{Nivel: Avanzado (Física Cuántica Contemporánea)}\\
    Considere una partícula unidimensional en un estado caracterizado por una función de onda con chirp o modulación de fase cuadrática:
    \begin{equation}
        \psi(x) = \left(\frac{1}{\pi\sigma^2}\right)^{1/4}\exp\left(-\frac{x^2}{2\sigma^2} + \iu\,\frac{\alpha}{2}x^2\right), \quad \alpha \in \mathbb{R}.
    \end{equation}
    \begin{enumerate}
        \item Demuestre que la incertidumbre espacial sigue siendo $\Delta x = \frac{\sigma}{\sqrt{2}}$.
        \item Demuestre que la incertidumbre en el momento se incrementa a:
        \begin{equation}
            \Delta p_x = \frac{\hbar}{\sqrt{2}\sigma}\sqrt{1 + \alpha^2\sigma^4}.
        \end{equation}
        \item Calcule la covarianza cuántica $\mathrm{Cov}(\hat{x}, \hat{p}_x) = \frac{1}{2}\ev{\{\hat{x}, \hat{p}_x\}}$ y demuestre que satura con igualdad exacta la relación generalizada de \textbf{Robertson--Schrödinger} \eqref{eq:robertson_schrodinger}, pero es estrictamente mayor que la cota simple de Heisenberg $\hbar/2$.
    \end{enumerate}
\end{ejerciciobox}

\begin{figure}[H]
    \centering
    % Fig. 2: gaussiana con chirp (sigma = 1, hbar = 1)
\begin{tikzpicture}
\begin{groupplot}[
  group style={group size=2 by 1, horizontal sep=2.2cm},
  width=0.5\linewidth, height=5.2cm,
  axis lines=left, every axis plot/.append style={thick}]
\nextgroupplot[domain=-3.5:3.5, samples=400, xlabel={$x/\sigma$}, ylabel={$\mathrm{Re}\,\psi(x)$},
  title={\small Parte real de $\psi$}, ymin=-0.85, ymax=0.85,
  legend style={font=\scriptsize, draw=none, fill=none, at={(0.98,0.98)}, anchor=north east}]
  \addplot[uni-blue]{exp(-x^2/2)*cos(deg(0*x^2/2))*0.7511};
  \addplot[alert-border]{exp(-x^2/2)*cos(deg(3*x^2/2))*0.7511};
  \addplot[gray, dashed, thin, forget plot]{exp(-x^2/2)*0.7511};
  \addplot[gray, dashed, thin, forget plot]{-exp(-x^2/2)*0.7511};
  \legend{$\alpha\sigma^2=0$, $\alpha\sigma^2=3$}
\nextgroupplot[xlabel={$x/\sigma$}, ylabel={$\sigma \cdot p/\hbar$},
  title={\small Elipse de incertidumbre en $(x,p)$},
  xmin=-1.3, xmax=1.3, ymin=-3.2, ymax=3.2, axis lines=left,
  legend style={font=\scriptsize, draw=none, fill=none, at={(0.02,0.98)}, anchor=north west}]
  \addplot[uni-blue, domain=0:360, samples=100]
    ({0.7071*cos(x)}, {0.7071*sin(x)});
  \addplot[alert-border, domain=0:360, samples=100]
    ({0.7071*cos(x)}, {1.4142*cos(x)+0.7071*sin(x)});
  \legend{$\alpha\sigma^2=0$, $\alpha\sigma^2=2$}
\end{groupplot}
\end{tikzpicture}

    \caption{Gaussiana con chirp (Ejercicio Propuesto 3). Izquierda: la envolvente es la misma, pero las oscilaciones se aceleran en los bordes. Derecha: la elipse de incertidumbre en el plano $(x,p)$ se inclina (covarianza $\hbar\alpha\sigma^2/2\neq0$) y $\Delta p_x$ crece, mientras $\Delta x$ no cambia.}
    \label{fig:chirp}
\end{figure}

\subsubsection[Ejercicio 4: Acoplamiento $LS$ y centro de gravedad]{Ejercicio Propuesto 4: Acoplamiento Espín-Órbita $\hat{H}_{SO}$ y Regla del Centro de Gravedad}
\begin{ejerciciobox}[Ejercicio Propuesto 4: Estructura Fina y Desdoblamiento Espín-Órbita]
    \textbf{Nivel: Avanzado}\\
    En la estructura fina de un átomo hidrogenoide, el acoplamiento magnético entre el espín del electrón ($s = 1/2$) y su movimiento orbital ($\ell = 1$, nivel $2p$) genera el término de interacción:
    \begin{equation}
        \hat{H}_{SO} = \xi(\hat{\vec{L}}\cdot\hat{\vec{S}}), \quad \xi > 0,
    \end{equation}
    donde la constante $\xi \equiv \ev{\xi(r)}_{2p}$ es el valor esperado radial de la función de acoplamiento espín-órbita $\xi(r)$ evaluada en el subespacio del nivel $2p$ (consistente con la Sección~\ref{subsec:acoplamiento_espin_orbita}).
    \begin{enumerate}
        \item A partir de la definición del momento angular total $\hat{\vec{J}} = \hat{\vec{L}} + \hat{\vec{S}}$, demuestre la identidad operacional:
        \begin{equation}
            \hat{\vec{L}}\cdot\hat{\vec{S}} = \frac{1}{2}\left( \hat{J}^2 - \hat{L}^2 - \hat{S}^2 \right).
        \end{equation}
        \item Aplicando las reglas de adición de momentos angulares ($|\ell - s| \le j \le \ell + s$), determine los dos valores permitidos de $j$, su notación espectroscópica estándar ($^2P_{1/2}$ y $^2P_{3/2}$) y sus respectivas degeneraciones $g_j = 2j + 1$.
        \item Calcule el corrimiento de energía $\Delta E(j)$ inducido por $\hat{H}_{SO}$ sobre cada uno de los dos multipletes y demuestre que la separación energética (desdoblamiento de estructura fina) es:
        \begin{equation}
            \Delta E_{\text{fina}} = E(^2P_{3/2}) - E(^2P_{1/2}) = \frac{3}{2}\xi\hbar^2.
        \end{equation}
        \item Demuestre que se satisface la \textbf{regla de suma del centro de gravedad}:
        \begin{equation}
            \sum_j (2j + 1)\Delta E(j) = 0,
        \end{equation}
        y explique su significado físico en términos de la invariancia de la traza del operador perturbativo.
    \end{enumerate}
\end{ejerciciobox}

\begin{figure}[H]
\centering
\begin{tikzpicture}[>=Stealth,font=\small]
  \def\sc{2.0}   % cm por unidad de  xi*hbar^2

  % eje de energía
  \draw[->,grisUNI] (-0.8,-2.9)--(-0.8,2.0) node[above]{$E$};

  % nivel sin acoplamiento
  \draw[very thick] (0,0)--(2.2,0);
  \node[anchor=east] at (-0.05,0) {$2p$};
  \node[above,font=\footnotesize] at (1.1,0.05) {$g=6$};

  % niveles desdoblados
  \draw[very thick,rojoUNI] (5,{0.5*\sc})--(7.2,{0.5*\sc});
  \draw[very thick,azulUNI] (5,{-1.0*\sc})--(7.2,{-1.0*\sc});
  \draw[dotted,grisUNI] (2.2,0)--(5,{0.5*\sc});
  \draw[dotted,grisUNI] (2.2,0)--(5,{-1.0*\sc});

  \node[rojoUNI,anchor=south west,font=\small] at (5,{0.5*\sc+0.05})
       {${}^{2}\!P_{3/2}$\; ($g=4$)};
  \node[rojoUNI,anchor=north west,font=\footnotesize] at (5,{0.5*\sc-0.05})
       {$\Delta E=+\tfrac12\,\xi\hbar^2$};
  \node[azulUNI,anchor=south west,font=\small] at (5,{-1.0*\sc+0.05})
       {${}^{2}\!P_{1/2}$\; ($g=2$)};
  \node[azulUNI,anchor=north west,font=\footnotesize] at (5,{-1.0*\sc-0.05})
       {$\Delta E=-\,\xi\hbar^2$};

  % centro de gravedad
  \draw[dashed,naranjaUNI] (2.2,0)--(7.5,0);
  \node[naranjaUNI,anchor=south west,font=\footnotesize,align=left] at (7.7,0.1)
       {centro de gravedad:\\ $4\!\cdot\!\tfrac12+2\!\cdot\!(-1)=0$};

  % desdoblamiento
  \draw[<->,thick] (7.5,{-1.0*\sc})--(7.5,{0.5*\sc});
  \node[anchor=west,font=\footnotesize,align=left] at (7.7,{-0.75*\sc})
       {$\Delta E_{\rm fina}=\tfrac32\,\xi\hbar^2$};
\end{tikzpicture}
\caption{Estructura fina del nivel $2p$ ($\ell=1$, $s=\tfrac12$) por
$\hat H_{SO}=\xi\,\hat{\vec L}\cdot\hat{\vec S}$ con $\xi>0$. La perturbación
reparte el nivel en $^2P_{3/2}$ (subida) y $^2P_{1/2}$ (bajada) de modo que
$\sum_j(2j+1)\Delta E_j=0$: la traza de $\hat H_{SO}$ en el subespacio de
dimensión 6 es nula y la energía media no cambia.}
\label{fig:estructura-fina}
\end{figure}


\subsubsection[Ejercicio 5: Estados de Bell y traza parcial]{Ejercicio Propuesto 5: Producto Tensorial, Estados de Bell y Traza Parcial}
\begin{ejerciciobox}[Ejercicio Propuesto 5: Entrelazamiento Cuántico y Traza Parcial]
    \textbf{Nivel: Avanzado}\\
    Sean $\mathcal{H}_A$ y $\mathcal{H}_B$ dos espacios de Hilbert bidimensionales (qubits) con bases ortonormales $\{\ket{0}, \ket{1}\}$.
    \begin{enumerate}
        \item Demuestre formalmente por contradicción que el \textbf{estado de Bell}:
        \begin{equation}
            \ket{\Phi^+} \equiv \frac{1}{\sqrt{2}}\left( \ket{0}_A \otimes \ket{0}_B + \ket{1}_A \otimes \ket{1}_B \right)
        \end{equation}
        no puede factorizarse como un producto tensorial simple $\ket{\psi_A} \otimes \ket{\psi_B}$, demostrando que es un estado cuántico genuinamente entrelazado.
        \item Construya el operador de densidad global del estado puro $\hat{\rho} \equiv \ket{\Phi^+}\bra{\Phi^+}$ en forma matricial de $4\times 4$ en la base ordenada $\{\ket{00}, \ket{01}, \ket{10}, \ket{11}\}$.
        \item Calcule la matriz de densidad reducida del subsistema $A$ mediante la \textbf{traza parcial} sobre el subsistema $B$:
        \begin{equation}
            \hat{\rho}_A \equiv \mathrm{Tr}_B(\hat{\rho}) = \sum_{k=0}^1 (\hat{I}_A \otimes \bra{k}_B)\hat{\rho}(\hat{I}_A \otimes \ket{k}_B).
        \end{equation}
        Verifique que $\hat{\rho}_A = \frac{1}{2}\hat{I}_2$ y concluya por qué el subsistema individual $A$ se encuentra en un \textbf{estado maximalmente mezclado} (entropía cuántica máxima de von Neumann), a pesar de que el estado global del par es estrictamente puro.
    \end{enumerate}
\end{ejerciciobox}
